% =====================================================================
\chapter{Valor absoluto y desigualdades}
\label{tema:2}
% =====================================================================

Gran parte del trabajo matemático en aprendizaje automático consiste en \emph{acotar}: garantizar que un error es pequeño, que una pérdida no es negativa o que un método converge. Este tema reúne las desigualdades elementales que se usan para ello. El valor absoluto y la desigualdad triangular siguen el capítulo~1 de Spivak~\cite{spivak}; las inecuaciones, el §2.7 de OpenStax~\cite{openstax}; las demostraciones por inducción, el capítulo~10 de Hammack~\cite{hammack}.

\section{Valor absoluto}

\begin{definicion}[Valor absoluto]\label{def:abs}
Para $x\in\R$ se define $\abs{x}=x$ si $x\ge 0$ y $\abs{x}=-x$ si $x<0$. La \emph{distancia} entre $x$ e $y$ es $\abs{x-y}$.
\end{definicion}

\begin{proposicion}\label{prop:abs}
Para todos $x,y\in\R$ y $r\ge 0$:
\begin{enumerate}
  \item $\abs{x}\ge 0$, con igualdad si y solo si $x=0$; $\abs{-x}=\abs{x}$;
  \item $\abs{x}=\max(x,-x)$ y $-\abs{x}\le x\le\abs{x}$;
  \item $\abs{x}=\sqrt{x^2}$ y $\abs{xy}=\abs{x}\abs{y}$;
  \item $\abs{x}\le r$ si y solo si $-r\le x\le r$; $\abs{x}\ge r$ si y solo si $x\le -r$ o $x\ge r$.
\end{enumerate}
\end{proposicion}

\section{La desigualdad triangular}

\begin{teorema}[Desigualdad triangular]\label{teo:triangular}
Para todos $x,y\in\R$, $\abs{x+y}\le\abs{x}+\abs{y}$, con igualdad si y solo si $xy\ge0$.
\end{teorema}

\begin{corolario}[Desigualdad triangular inversa]\label{cor:triangular-inversa}
Para todos $x,y\in\R$, $\bigl|\abs{x}-\abs{y}\bigr|\le\abs{x-y}$.
\end{corolario}

La desigualdad triangular se generaliza a $n$ sumandos, $\abs{\sum_i x_i}\le\sum_i\abs{x_i}$ (\cref{eje:2-triangular-n}), y a vectores, donde $\abs{\cdot}$ se sustituye por una norma (módulo~2).

\section{Inecuaciones con valor absoluto}

Por la \cref{prop:abs}(iv), la condición $\abs{x-a}<r$ describe el intervalo $(a-r,a+r)$ de centro $a$ y radio $r$. Las inecuaciones con varios valores absolutos se resuelven separando casos según el signo de cada expresión.

\addtocounter{definicion}{1}%

\section{Desigualdad entre las medias aritmética y geométrica}

\begin{proposicion}[Caso de dos números]\label{prop:amgm2}
Si $a,b\ge0$, entonces $\sqrt{ab}\le\dfrac{a+b}{2}$, con igualdad si y solo si $a=b$.
\end{proposicion}

\begin{lema}\label{lem:producto1}
Si $x_1,\dots,x_n>0$ y $x_1x_2\cdots x_n=1$, entonces $x_1+\dots+x_n\ge n$, con igualdad si y solo si todos valen~1.
\end{lema}

\begin{teorema}[Desigualdad de las medias]\label{teo:amgm}
Si $a_1,\dots,a_n>0$, entonces
\[
\sqrt[n]{a_1a_2\cdots a_n}\le\frac{a_1+a_2+\dots+a_n}{n},
\]
con igualdad si y solo si $a_1=a_2=\dots=a_n$.
\end{teorema}

\addtocounter{definicion}{1}%

\section{Desigualdades de Bernoulli y de Cauchy--Schwarz}

\begin{teorema}[Desigualdad de Bernoulli]\label{teo:bernoulli}
Si $x\ge-1$ y $n\in\N$, entonces $(1+x)^n\ge1+nx$.
\end{teorema}

\begin{teorema}[Desigualdad de Cauchy--Schwarz]\label{teo:cs}
Para reales $a_1,\dots,a_n$ y $b_1,\dots,b_n$,
\[
\Bigl(\sum_{i=1}^na_ib_i\Bigr)^2\le\Bigl(\sum_{i=1}^na_i^2\Bigr)\Bigl(\sum_{i=1}^nb_i^2\Bigr).
\]
\end{teorema}

\addtocounter{definicion}{1}%

\addtocounter{definicion}{1}%

